\subsection{An all-order obstruction for two minimal rows}

\begin{lemma}[Finite affine acts on a line]\label{C:finite-line-act-cups}
Let $M$ be a monoid and let $T\subset\mathbb Q$ be a nonempty finite
right $M$-set whose action maps have the form
\[
 t\longmapsto\chi_m t+\beta_m,
 \qquad \chi_m\in\{0,1,-1\}.
\]
Every positive-degree rational cohomology product of the act-cone
category $\mathcal C(M,T)$ is zero.
\end{lemma}
\begin{proof}
For a singleton $T$, the edge bar is $BM$ and its projection to
$BM$ is the identity, so the cone is contractible.  Suppose $T$
has at least two elements, and put $I=[\min T,\max T]$.
An affine map of slope $1$ preserving a finite subset of the line
must be the identity: a nonzero translation cannot permute that
finite set.  A map of slope $-1$ is a bijection of $T$, so its
translation term is necessarily $\min T+\max T$.  Thus all such
maps are the same reflection of $I$.  The remaining maps are
constants at points of $T$.  The given action consequently extends
to an actual affine right action on the contractible interval $I$.
Its multiplication identities hold on $I$, since they hold as
affine identities on the at least two points of $T$.

The levelwise weak equivalence from the interval bar to $BM$ gives
\[
 B\mathcal C(M,T)
 \simeq\operatorname{hocofib}\bigl(B(T,M,*)\longrightarrow BM\bigr)
 \simeq\operatorname{hocofib}\bigl(B(T,M,*)
                                  \longrightarrow B(I,M,*)\bigr).
                                                               \tag{\ref{C:finite-line-act-cups}.1}
\]
The last cofiber is the pointed bar with fiber $I/T$, with its
basepoint bar collapsed.  This follows level by level from the
cofibration $T\subset I$; realization preserves this cofiber.
Writing the elements of $T$ in increasing order identifies $I/T$
with a wedge of $|T|-1$ circles, one for each consecutive interval.
Each action map is the identity, the common reflection, or the
constant pointed map.  The reflection permutes these circle
summands and reverses their orientations.

For an integer $d\geq2$, take the degree-$d$ circle map on every
summand, using its linear interval parameter modulo its endpoints.
It commutes strictly with the reflection: on the circle parameter,
the two expressions are $d(1-t)$ and $1-dt$, equal modulo one.
It also commutes with every constant map and the identity.
It therefore induces a self-map $F_d$ of the pointed bar in
(\ref{C:finite-line-act-cups}.1).

The reduced cellular chain complex of $I/T$ is concentrated in
degree one, with a free generator for each consecutive interval.
All action maps are cellular, and $F_d$ is multiplication by $d$
on this entire reduced chain complex.  The normalized pointed-bar
cellular complex thus has the same literal multiplication-by-$d$
map in every degree.  Hence $F_d^*$ is multiplication by $d$ on
every positive-degree rational cohomology group.  For positive
classes $u,v$, naturality of the cup product gives
\[
 d\,uv=F_d^*(uv)=F_d^*(u)F_d^*(v)=d^2uv.
\]
Since $d^2-d\ne0$ in $\mathbb Q$, this proves $uv=0$.
\end{proof}

\begin{theorem}[Two minimal rows exclude every act-cone square]
\label{C:two-row-act-cups}
Let $M$ be a finite monoid whose minimal Rees matrix ideal has at
most two rows.  For every nonempty finite right $M$-set $S$, all
products
\[
 H^2(B\mathcal C(M,S);\mathbb Q)\otimes
 H^2(B\mathcal C(M,S);\mathbb Q)
 \longrightarrow H^4(B\mathcal C(M,S);\mathbb Q)
\]
are zero.  In particular no such act-cone category has the homotopy
type of $\mathbb{CP}^2$.  There is no restriction on the order of
$M$, the size of $S$, or the maximal group of the minimal ideal.
The opposite assertion holds for left acts and at most two minimal
columns.
\end{theorem}
\begin{proof}
Put $R=\mathbb QM$.  Choose a minimal idempotent $e$ and write
$H=eMe$ for its maximal group.  The averaging element
\[
 \bar e=|H|^{-1}\sum_{h\in H}h
\]
is an idempotent of $R$.  The projective left module $P=R\bar e$
has one basis vector for each row of the minimal ideal: $Me$ is
the disjoint union of the corresponding free right $H$-orbits,
and multiplying by $\bar e$ averages each orbit.  Left
multiplication by every $m\in M$ sends each such basis vector to
another row basis vector.  Its augmentation $P\to\mathbb Q$
therefore sends each basis vector to $1$.

If there is only one row, this augmentation is an isomorphism.
The trivial module is projective, and
Lemma~\ref{C:act-cone-Tor} gives $H^2(B\mathcal C(M,S);\mathbb Q)=0$.
Suppose there are two rows.  The kernel of the augmentation is
the one-dimensional row-difference module, giving an exact sequence
\[
 0\longrightarrow\mathbb Q_\chi\longrightarrow P
      \longrightarrow\mathbb Q\longrightarrow0,
 \qquad \chi:M\longrightarrow\{0,1,-1\}.             \tag{\ref{C:two-row-act-cups}.1}
\]
Here $\chi_m$ is the action of the two-point row transformation on
its augmentation line: it is $1$ for the identity, $-1$ for the
interchange, and $0$ for a constant.  In particular $\chi$ is an
actual multiplicative character, including its zero values.

Let $L_S=\ker(\mathbb QS\to\mathbb Q)$, a right $R$-module,
and write $D L_S=\operatorname{Hom}_{\mathbb Q}(L_S,\mathbb Q)$
with the contragredient left action.  The natural act-cone formula
and finite-dimensional vector-space duality give
\[
 H^2(B\mathcal C(M,S);\mathbb Q)
 \cong D\operatorname{Tor}_1^R(L_S,\mathbb Q)
 \cong\operatorname{Ext}_R^1(\mathbb Q,D L_S).        \tag{\ref{C:two-row-act-cups}.2}
\]
The second isomorphism follows by dualizing the complex obtained
from a projective resolution of the left module $\mathbb Q$;
it is natural in every map of right acts $S\to T$.
Applying $\operatorname{Hom}_R(-,D L_S)$ to
(\ref{C:two-row-act-cups}.1), and using projectivity of $P$, gives
a surjection
\[
 \operatorname{Hom}_R(\mathbb Q_\chi,D L_S)
 \longrightarrow H^2(B\mathcal C(M,S);\mathbb Q).     \tag{\ref{C:two-row-act-cups}.3}
\]
Thus every degree-two class has a representative consisting of a
linear functional $\ell:L_S\to\mathbb Q$ satisfying
\[
 \ell(v m)=\chi_m\ell(v)\qquad(v\in L_S, m\in M).  \tag{\ref{C:two-row-act-cups}.4}
\]
No assertion that $L_S$ itself is annihilated by the minimal ideal
is required.

Fix $s_0\in S$ and define
\[
 q(s)=\ell([s]-[s_0]),\qquad T=q(S)\subset\mathbb Q.
\]
For every $s\in S$ and $m\in M$, equation
(\ref{C:two-row-act-cups}.4) gives
\[
 q(sm)=\chi_m q(s)+q(s_0m).                           \tag{\ref{C:two-row-act-cups}.5}
\]
Consequently $q:S\to T$ is a map of finite right $M$-sets when
$T$ is given these affine action maps.  This is well defined even
when $q$ is not injective.  The functional
\[
 \ell_T\left(\sum_{t\in T}a_t[t]\right)
       =\sum_{t\in T}a_t t,
 \qquad \sum_{t\in T}a_t=0,
\]
is $\chi$-equivariant by (\ref{C:two-row-act-cups}.5), since the
translation term is multiplied by $\sum a_t=0$.
Moreover $\ell=\ell_T\circ L(q)$.

Let $x$ be the original degree-two class, and let $y$ be the
class for $\mathcal C(M,T)$ obtained from $\ell_T$ by the same
connecting homomorphism (\ref{C:two-row-act-cups}.3).
Naturality of that connecting map and of
(\ref{C:two-row-act-cups}.2) proves
\[
 x=(B\mathcal C(M,q))^*y.
\]
The target act is affine on a finite rational subset of a line.
Lemma~\ref{C:finite-line-act-cups} gives $y^2=0$, and hence $x^2=0$.
This holds for every $x$.  Finally rational polarization,
$(x+z)^2-x^2-z^2=2xz$, proves that every product of two degree-two
classes is zero.  Reversing the category interchanges right acts
with left acts and rows with columns, proving the opposite statement.
\end{proof}

\begin{corollary}\label{C:two-row-biset-cups}
The same exclusion applies to a triangular biset category whose
first endomorphism monoid has at most two minimal rows and whose
other endomorphism monoid $N$ has projective trivial rational module
on the side used for its coinvariants.
\end{corollary}
\begin{proof}
Take the actual finite orbit act $S/N$.  Its quotient map from the
cross-biset is equivariant for the first monoid and defines a functor
from the biset category to $\mathcal C(M,S/N)$.  In every degree of
the first monoid's bar, exact $N$-coinvariants identify the second
bar rationally with the orbit set.  Thus the map of edge bars is a
rational homology equivalence.  The vertex map $BN\to *$ is also
a rational homology equivalence.  The mixed-nerve homotopy-pushout
comparison gives a rational homology equivalence of the complete
classifying spaces, and therefore an isomorphism of their rational
cohomology rings.  Theorem~\ref{C:two-row-act-cups} applies to the
target.  This is an actual quotient-functor comparison; no finite-type
property of an individual vertex or edge space is inferred.
\end{proof}
